%% An example of the article based on superfri.cls class file of
%% ``Supercomputing Frontiers and Innovations. An International Journal''
%% http://superfri.org/.

\documentclass{superfri}
\usepackage{graphicx}
\usepackage[hidelinks]{hyperref}

% ------------
\bibliographystyle{superfri}

\begin{document}

\author{Alice~M.~Scientist\footnote{\label{susu}South Ural State University, Chelyabinsk, Russian Federation}$^{,2}$\orcidID{0000-1111-2222-3333} \and  Bob~A.~Scientist\footnoteref{susu} \and Derek~R.~Author\footnote{\label{msu}Lomonosov Moscow State University, Moscow, Russian Federation}\orcidID{1111-2222-3333-4444}
}

\title{SuperFrI Article Template\thanks{The paper is recommended for publication by the Program Committee of the International Scientific Conference ``Parallel computational technologies (PCT) 2018''.}}

\maketitle{}

\begin{abstract}%
This template is devoted to helping you format your article the proper way in
\LaTeX.

\keywords{superfri, template, article formatting}
\end{abstract}

% -----------------------------------------------------------------------
\section*{Introduction}
\label{sec:intro}
The template demonstrates~\cite{DBLP:conf/ics/2015} how to use the \emph{superfri} class. It could be
used as a basis~\cite{DBLP:conf/iccS/AltintasNLKDS16} for your article and meets all the formatting requirements of
the SuperFri journal. Consult the author
guidelines for details on that matter.

\section{The Class}
\label{sec:class}

\subsection{Installation}
\label{sec:install}
The class consists of a single file \verb=superfri.cls= which should reside
along with your manuscript file.

\subsection{Features}
\label{sec:features}
The class is based on the standard~\cite{DBLP:books/sp/byeon2010/RiedelSMWL10} article class and supports all of its
features, as well as these:
\begin{itemize}
\item title and author formatting;
\item abstract and keywords;
\item theorems, definitions, and proofs.
\end{itemize}

Begin your paper~\cite{DBLP:series/cogtech/354089408} with the title and authors' names. After filling the
necessary info in with \verb=\title= and \verb=\author= commands put it all
together~\cite{SIV2008} using the \verb=\maketitle= command.

Define the structure of your article using up to three levels of nesting with
the \verb=\section=, \verb=\subsection=, and \verb=\subsubsection= commands.
You can use the environments \verb=theorem=, \verb=proof=, \verb=lemma=, etc.

For further directions read the comments in \verb=superfri.cls= as a manual~\cite{DBLP:journals/pcs/SokolinskyS16}.

\subsubsection{Subsubsection title}

The equations look like
\begin{equation}
a^2 + b^2 = c^2 + \int_1^2 x+y \, \mathrm{d}y,
\end{equation}
where $2$ is a number.

An example of a simple figure~\cite{DBLP:journals/corr/Zhang16b} is shown in~\figref{pic/pic}, while a table
is in~\tabref{tab-mytable}.

\fig{width=5cm}{pic/pic}{An example of a simple figure}

Figure~\ref{fig:fig} demonstrates formatting of a complex figure.

\begin{figure}[!ht]
	\begin{minipage}{.5\linewidth}
		\centering
		\includegraphics[width=0.65\linewidth]{pic/pic}
		\subcaption{This is subcaption for A subfigure}\label{subfig:a}
	\end{minipage}
	\begin{minipage}{.5\linewidth}
		\centering
		\includegraphics[width=0.65\linewidth]{pic/pic}
		\subcaption{This is subcaption for B subfigure}\label{subfig:b}
	\end{minipage}
	\vspace{3pt}
	\caption{An example of a complex figure}
	\label{fig:fig}
\end{figure}


\tab{tab-mytable}{An example of a table}{
	\begin{tabular}{rcl}
	\hline
	one & two & three \\
	\hline
	a & b & c \\
	$\alpha$ & $\beta$ & $\gamma$ \\
	$x$ & $y$ & $z$ \\
	\hline
	\end{tabular}
}

\section*{Acknowledgements}
Put the acknowledgements after the last section, like this.

\openaccess

\bibliography{example}
\end{document}
